En la siguiente figura se muestra un triángulo $ABC$.

\begin{center}
\begin{tikzpicture}[scale=0.8]

% Triángulo
\coordinate (A) at (0,0);
\coordinate (B) at (6,0);
\coordinate (C) at (2,4);

\draw (A)--(B)--(C)--cycle;

% Etiquetas
\node[below left] at (A) {A};
\node[below right] at (B) {B};
\node[above] at (C) {C};

\end{tikzpicture}
\end{center}

Para construir la bisectriz del ángulo en el vértice $A$, siga los siguientes pasos utilizando regla y compás:

\begin{enumerate}
\item Con centro en $A$, trace un arco que corte los lados $AB$ y $AC$.
\item Sin cambiar la abertura del compás, trace arcos con centro en cada uno de los puntos de intersección obtenidos en los lados.
\item Marque el punto donde se intersectan estos dos arcos.
\item Trace una recta desde el vértice $A$ hasta el punto de intersección de los arcos.
\end{enumerate}

Realice la construcción sobre el triángulo y explique con sus palabras que es la bisectriz del ángulo.